% ===================================================================
%  तक्ता क्रमांक ५ — घर आणि त्याचे बांधकाम.
%
%  Traduction, faite en 2026, en marathi d'aujourd'hui, sur l'ido de
%  texto/io. Regles de la colonne : voir 10-takta-01.tex. Le tableau
%  est un DIALOGUE, et les deux volumes composent le nom du parleur
%  en petites capitales : on garde \textsc{}, que le rendu laisse
%  tomber en gardant le texte. Les renvois du plan sont des LETTRES
%  — (a) le couloir, (b) le salon... — et l'on ecrit « (1) » pour la
%  salle de bain la ou l'imprimeur l'a compose ainsi au lieu de (l).
%
%  LE MARATHI REPOND A LA NOTE (*), ET C'EST LE PLUS BEL ECHO DE LA
%  COLONNE. La note occupe six lignes a expliquer que l'ido n'a pas
%  de racine pour la « solive » et qu'il faudrait en adopter une :
%  « balk-o, radiko teknikala til nun ne adoptita ». Le marathi a le
%  mot depuis toujours et sans y penser : वासा, वासे — la piece de
%  bois equarrie qui porte le plancher et le plafond, ni तुळई (la
%  poutre) ni लहान तुळई. Six langues cherchaient un mot ; la
%  septieme le disait tous les jours.
%
%  UN NOM DE METIER PAR METIER, et le plus souvent le nom d'une
%  caste d'artisans, que le Maharashtra porte encore comme nom de
%  famille :
%
%    गवंडी (108)      le macon
%    पाथरवट (83)      le tailleur de pierre
%    सुतार (115)      le charpentier
%    लोहार (125)      le forgeron
%    माळी (129)       le jardinier
%    रंगारी (102)     le peintre
%    मोतद्दार (133)    le palefrenier — du persan, par les Marathes
%    बिगारी (146)     le manœuvre
%    मातीकाम करणारा (81) le terrassier
%
%  ET LA COLONNE A UN TROU, LE DEUXIEME DE LA SERIE. L'ido separe le
%  karpentisto (qui fait la charpente) du menuzisto (qui fait les
%  portes et les parquets) ; le marathi dit सुतार pour les deux et
%  distingue par le travail, non par le nom. On ne fabrique pas un
%  mot pour boucher un trou : on ecrit सुतार deux fois, comme le
%  ferait un Marathe.
%
%  LE POINT, ET NON LE DANDA. Le hindi de la colonne voisine termine
%  ses phrases par « । » ; le marathi ecrit « . » depuis le XIXe
%  siecle. Deux colonnes en devanagari, deux ponctuations.
%
%  L'OUTILLAGE ET LA MAISON, mot pour mot :
%
%    ओळंबा (109)   le fil a plomb      ऐरण (127)    l'enclume
%    सांडशी (128)  les tenailles       कानस (98)    la lime
%    दंताळे (130)  le rateau           झारी (132)   l'arrosoir
%    कुदळ (131)    la beche            वाफे (150)   les planches
%    सज्जा (28)    le balcon           गच्ची (21)   la terrasse
%    फाटक (48)     le portail          अडसर (53)    le verrou
%    उंबरठा (10)   le seuil            चौथरा (9)    le perron
%    चर (46)       la tranchee         गारा (65)    le mortier
%
%  SEPT INVERSIONS DE RENVOI, toutes dues a la postposition : le
%  marathi place le possesseur AVANT le possede, et « X de Y » sort
%  donc a l'envers. Le remede est celui des tableaux 3 et 4 : nommer
%  d'abord ce que l'ido nomme d'abord, puis qualifier par une
%  seconde proposition, entre tirets.
% ===================================================================

%%K t05-apar-1 apar
\VUsaut{6.0mm}
\VUtitre{15.0pt}{60}{तक्ता क्रमांक ५}
\VUsaut{1.4mm}
\VUtitre{11.4pt}{0}{घर आणि त्याचे बांधकाम.}
\VUsaut{2.2mm}

%%K t05-tit-1 sub
\VUpk{10.2pt}{40}{चांगली भेट.}
\VUsaut{3.0mm}

%%K t05-01-1 p
\textsc{जॉन}. --- काका, \VUgras{घर} कसे बांधतात हे मला फार जाणून घ्यायचे आहे.

%%K t05-01-2 p
\textsc{काका} \textsuperscript{(80)}. --- हे अगदी सोपे आहे रे मित्रा; माझ्याबरोबर चल, श्री.
बर्नार्ड \textsuperscript{(79)} जे घर बांधून घेत आहेत आणि ज्यात मी राहणार आहे, ते मी तुला
दाखवतो.

%%K t05-01-3 p
\textsc{जॉन}. --- मग या घराचा \VUgras{नकाशा} \textsuperscript{(76)} कोणी काढला?

%%K t05-01-4 p
\textsc{काका}. --- \VUgras{वास्तुविशारदाने} \textsuperscript{(75)}; त्यांनी फार स्पष्ट आराखडा
सादर केला, तो मंजूर झाला, आणि \VUgras{कंत्राटदाराने} \textsuperscript{(77)} कामाला सुरुवात
केली.

%%K t05-01-5 p
\textsc{जॉन}. --- आणि कंत्राटदार कोण आहेत?

%%K t05-01-6 p
\textsc{काका}. --- श्री. ब्लांको, तुझा मित्र जेकब याचे वडील. वास्तुविशारद श्री. रुफो
यांच्यासोबत ते मजुरांना कामाची दिशा देतात.

%%K t05-01-7 p
अरे! हे बघ, श्री. ब्लांको आपली \VUgras{मापपट्टी} \textsuperscript{(78)} घेऊन आणि श्री. रुफो
आपला नकाशा घेऊन अगदी वेळेवर येत आहेत.

%%K t05-01-8 p
नमस्कार, महाशय. कसे आहात?

%%K t05-01-9 p
\textsc{श्री. रुफो}. --- अगदी छान, धन्यवाद. नमस्कार जॉन, केवढा वाढला हा मुलगा!

%%K t05-01-10 p
\textsc{काका}. --- होय, खरेच, हा आता लहान माणूसच झाला आहे. फार गंभीर आहे; जे काही दिसते
त्याबद्दल त्याला सगळे सांगावे लागते. आज त्याला घर कसे बांधतात हे जाणून घ्यायचे आहे.

%%K t05-tit-2 sub
\VUpk{10.2pt}{40}{कामगार.}
\VUsaut{3.0mm}

%%K t05-01-11 p
\textsc{श्री. ब्लांको}. --- बरं, माझ्याबरोबर ये रे मित्रा, मी तुला लगेच समजावून सांगतो, कारण
शिकू पाहणारी मुले मला फार आवडतात.

%%K t05-01-12 p
हातात \VUgras{फावडे} \textsuperscript{(82)} घेतलेला तो मजूर दिसतो का? तो
\VUgras{मातीकाम करणारा} \textsuperscript{(81)} आहे. पाण्याचे नळ टाकण्यासाठी तो \VUgras{चर}
\textsuperscript{(46)} खणतो आहे. घर जिथे उभे राहणार तिथली जमीनही त्यानेच खणून काढली, म्हणजे
गवंड्याला चार \VUgras{भिंती} उभारता याव्यात. त्या भिंतींचा जमिनीखालचा भाग म्हणजे \VUgras{पाया}
\textsuperscript{(2)}. पायांमधल्या मोकळ्या जागेचे \VUgras{तळघर} \textsuperscript{(3)} होते.
त्या तळघराला \VUgras{तळखिडकी} \textsuperscript{(5)} नावाची एक लहान खिडकी, एक \VUgras{दार}
\textsuperscript{(4)} आणि एक जिना असतो.

%%K t05-01-13 p
\VUgras{गवंड्याने} \textsuperscript{(108)} \VUgras{दारांसाठी} \textsuperscript{(8)} आणि \VUgras{खिडक्यांसाठी} \textsuperscript{(17)} मोकळ्या जागा सोडत भिंती बांधत नेल्या.

%%K t05-01-14 p
\textsc{जॉन}. --- पण तो गवंडी तर मला दिसत नाही!

%%K t05-01-15 p
श्री. \textsc{ब्लांको}. --- त्याचे घरातले काम आता संपले आहे. तो \VUgras{तबेल्याजवळ}
\textsuperscript{(54)} आपल्या \VUgras{मचाणावर} \textsuperscript{(110)} आहे; तो
\VUgras{धुणीघराची} \textsuperscript{(56)} भिंत बांधतो आहे.

%%K t05-01-16 p
त्याच्याजवळ थापी, घमेले आणि \VUgras{ओळंबा} \textsuperscript{(109)} आहे. पण ही नावे तुझ्या
लक्षात राहणार नाहीत.

%%K t05-01-17 p
\textsc{जॉन}. --- नक्कीच राहतील, कारण मी लगेच काकांकडे लहान थापी आणि घमेले मागणार आहे, आणि
ओळंबा मी स्वतःच दोरीने तयार करीन.

%%K t05-tit-3 sub
\VUsaut{3.0mm}
\VUpk{10.2pt}{40}{साहित्य.}
\VUsaut{3.0mm}

%%K t05-01-18 p
\textsc{काका}. --- वा! फार छान. पण मला सांग जॉन, घर बांधायला काय काय लागते?

%%K t05-01-19 p
\textsc{जॉन}. --- \VUgras{चिरे} \textsuperscript{(59)} आणि \VUgras{गारा}
\textsuperscript{(65)} लागतात.

%%K t05-01-20 p
\textsc{काका}. --- बरोबर. आपल्या \VUgras{गाड्यावर} \textsuperscript{(136)} बसलेला तो मोठी टोपी
घातलेला माणूस बघ. तो \VUgras{गाडीवान} \textsuperscript{(135)} आहे. रोज तो इथे
\VUgras{दगडी ठोकळे} \textsuperscript{(60)} आणून टाकतो. तो आपली गाडी अंगणात रिकामी करतो, आणि
\VUgras{पाथरवट} \textsuperscript{(83)} ते ठोकळे घडवतात आणि आपल्या \VUgras{दगडी करवतीने}
\textsuperscript{(85)} कापतात. \VUgras{बिगारी} \textsuperscript{(146)} ते गवंड्याकडे नेतो, आणि
गवंडी ते जागच्या जागी बसवतो.

%%K t05-01-21 p
दरम्यान, \VUgras{गारा मळणारा} \textsuperscript{(87)} गारा तयार करतो. त्याला \VUgras{वाळू}
\textsuperscript{(62)}, \VUgras{चुना} \textsuperscript{(63)} आणि \VUgras{पाणी}
\textsuperscript{(64)} लागते. चुना त्याच्याजवळ \VUgras{पोत्यात} \textsuperscript{(71)} आहे,
\VUgras{गवंड्याचा हाताखालचा} \textsuperscript{(111)} त्याला \VUgras{ढकलगाडीवरून}
\textsuperscript{(112)} वाळू आणून देतो, आणि \VUgras{शिकाऊ पोऱ्या} \textsuperscript{(113)}
पंपाशी (58) आहे. तो \VUgras{बादली} \textsuperscript{(114)} भरतो. गारा लवकरच तयार होतो.
\VUgras{गारा वाहणारा} \textsuperscript{(107)} तो घेतो आणि शिडीवरून वर मचाणावर नेतो, जिथे गवंडी
घराची ती बाजू पुरी करतो आहे.

%%K t05-01-22 p
आणि आता, \VUgras{छप्पर} \textsuperscript{(31)} करणाऱ्या कामगाराला काय म्हणतात?

%%K t05-01-23 p
\textsc{जॉन}. --- तो \VUgras{छप्पर घालणारा} \textsuperscript{(118)} आहे.

%%K t05-tit-4 sub
\VUsaut{3.0mm}
\VUpk{10.2pt}{40}{घराचे छप्पर.}
\VUsaut{3.0mm}

%%K t05-01-24 p
श्री. \textsc{ब्लांको}. --- होय, पण त्याच्या आधी \VUgras{सुतार} \textsuperscript{(115)} येतो.

%%K t05-01-25 p
सुतार \VUgras{छपराचा सांगाडा} \textsuperscript{(35)} उभा करतो. तो \VUgras{तुळया}
\textsuperscript{(37)} \VUgras{मेखांनी} \textsuperscript{(117)} जोडतो. वर ते रुंद मखमली
बाह्यांचे कामगार दिसतात ना, ते सुतार आहेत. एक जण लहान तुळई धरून आहे आणि दुसरा आपल्या
\VUgras{कुऱ्हाडीने} \textsuperscript{(116)} मेख घडवतो आहे. \VUgras{धुराड्याजवळ}
\textsuperscript{(41)} जो आहे तो छप्पर घालणारा. तो (*) वाशांवर पट्ट्या ठोकतो, आणि पट्ट्यांवर
\VUgras{स्लेट} \textsuperscript{(66)}. कधी कधी तो कौलेही घालतो.

%%K t05-noto-1 noto
\VUnotes{143.2mm}{%
(*) \VUgras{Balk-o} = जर्मन \textit{Balken}, इंग्रजी \textit{joist}, फ्रेंच \textit{solive},
इटालियन \textit{travicello}, रशियन \textit{balka}, स्पॅनिश व पोर्तुगीज \textit{viga}. अजून
स्वीकारली न गेलेली, पण फ्रेंच \textit{solive} या शब्दाचा अर्थ देण्यासाठी सुचवण्याजोगी अशी एक
तांत्रिक धातू. ती \textit{trabo} (फ्रेंच \textit{poutre}) नाही, आणि लहान तुळईही नाही. ती
म्हणजे चौकोनी घडवलेली लाकडी दांडी, जी लाकडी जमीन आणि छत तोलून धरते.%
}

%%K t05-01-26 p
तबेल्याचे छप्पर \VUgras{कौलारू} \textsuperscript{(55)} आहे, घरांचे \VUgras{स्लेटी}
\textsuperscript{(32)} आणि व्हरांड्याचे \VUgras{जस्ती} \textsuperscript{(24)}.

%%K t05-01-27 p
\textsc{जॉन}. --- जस्ती छप्परेही छप्पर घालणाराच करतो का?

%%K t05-01-28 p
श्री. \textsc{ब्लांको}. --- नाही, ते \VUgras{जस्तकाम करणाऱ्याचे} \textsuperscript{(121)} किंवा
\VUgras{नळकारागिराचे} काम — घराच्या छपरावर \VUgras{छपरी खिडकी} \textsuperscript{(38)} बसवतो
आहे तोच. \VUgras{पावसाच्या नळ्याही} \textsuperscript{(43)} आणि \VUgras{छपराच्या पन्हाळीही}
\textsuperscript{(42)} तोच बसवतो. जस्त आणि पत्रा तो डाग देऊन जोडतो. \VUgras{तडितरक्षक}
\textsuperscript{(40)} आणि \VUgras{वाऱ्याचा बाण} \textsuperscript{(39)} त्यानेच
\VUgras{छपराच्या टोकावर} \textsuperscript{(36)} बसवले.

%%K t05-01-29 p
\textsc{जॉन}. --- मग घर पुरे झाले का?

%%K t05-tit-5 sub
\VUsaut{3.0mm}
\VUpk{10.2pt}{40}{हत्यारे.}
\VUsaut{3.0mm}

%%K t05-01-30 p
श्री. \textsc{ब्लांको}. --- पूर्ण नाही. \VUgras{गिलावा करणारा} \textsuperscript{(99)}
\VUgras{विटांनी} \textsuperscript{(61)} त्या भिंती बांधतो, ज्या घराचे वेगवेगळ्या खोल्यांत
विभाजन करतात. \VUgras{गिलाव्याने} \textsuperscript{(70)} तो \VUgras{छतांना}
\textsuperscript{(26)} आणि भिंतींना गुळगुळीत करतो. आता तो \VUgras{व्हरांड्याचे}
\textsuperscript{(33)} छत करतो आहे. गवंड्याप्रमाणे त्याच्याजवळही \VUgras{थापी}
\textsuperscript{(100)} आणि \VUgras{घमेले} \textsuperscript{(101)} आहे.

%%K t05-01-31 p
नंतर सुतार दारे, खिडक्या, \VUgras{लाकडी फरशी} \textsuperscript{(16)} आणि \VUgras{झडपा}
\textsuperscript{(19)} करतो.

%%K t05-01-32 p
आता तो अंगणात आपल्या \VUgras{कामाच्या बाकड्यावर} \textsuperscript{(90)} काम करतो आहे.
त्याच्याजवळ \VUgras{करवत} \textsuperscript{(94)} आहे — लाकूड (69) कापण्यासाठी —, \VUgras{रंधा}
\textsuperscript{(91)}, \VUgras{पकड} \textsuperscript{(92)}, \VUgras{छिन्नी}
\textsuperscript{(94 bis)}, \VUgras{कंपास} \textsuperscript{(95)} आणि लाकडी \VUgras{हातोडा}
\textsuperscript{(95 bis)}.

%%K t05-01-33 p
\textsc{जॉन}. --- वा! पण गवंडीकामापेक्षा सुतारकाम जास्त मजेदार आहे. काका, तू मला कामाचे बाकडे,
करवत आणि रंधा घेऊन देशील ना? कारण मी लवकरच मेदोरसाठी एक लहानसे घर तयार करणार आहे.

%%K t05-01-34 p
\textsc{काका}. --- हो रे बाळा.

%%K t05-01-35 p
\textsc{जॉन}. --- मला किती आनंद झाला आहे! आणि प्रवेशद्वाराजवळ उभा आहे तो कोण?

%%K t05-tit-6 sub
\VUsaut{3.0mm}
\VUpk{10.2pt}{40}{रंगकाम.}
\VUsaut{3.0mm}

%%K t05-01-36 p
श्री. \textsc{ब्लांको}. --- तो \VUgras{रंगारी} \textsuperscript{(102)} आहे. \VUgras{रंगाचे}
\textsuperscript{(103)} भांडे आणि आपला \VUgras{कुंचला} \textsuperscript{(104)} घेऊन तो आपल्या
शिडीवर उभा आहे. प्रवेशद्वाराचे लाकूड जास्त काळ टिकावे म्हणून तो त्याला रंग देतो.

%%K t05-01-37 p
तोच खोल्यांमध्ये कागदही चिकटवतो.

%%K t05-01-38 p
\VUgras{सजावटकार} \textsuperscript{(107)} \VUgras{शिडीवर} \textsuperscript{(106)} आहे, \VUgras{सज्ज्यावर} \textsuperscript{(28)}, आणि \VUgras{चिक} \textsuperscript{(29)} बसवतो आहे.

%%K t05-01-39 p
मोठ्या खिडकीजवळ उभा असलेला कामगार \VUgras{काचकाम करणारा} \textsuperscript{(96)} आहे. तो
\VUgras{तावदाने} \textsuperscript{(18)} \VUgras{पुट्टीने} \textsuperscript{(73)} बसवतो.
तळघराच्या दाराजवळ गुडघे टेकून बसलेला तो दुसरा कामगार \VUgras{कुलूपकाम करणारा}
\textsuperscript{(97)} आहे. तो \VUgras{कुलूप} \textsuperscript{(6)} बसवतो आहे. त्याची
\VUgras{कानस} \textsuperscript{(98)} जवळच आहे.

%%K t05-01-40 p
\textsc{जॉन}. --- लोखंडावर आपल्या \VUgras{हातोड्याने} \textsuperscript{(84)} घाव घालणारा तो
कामगारही कुलूपकाम करणाराच आहे का?

%%K t05-01-41 p
श्री. \textsc{ब्लांको}. --- अधिक नेमके सांगायचे तर तो लोहार (125) आहे. तो
\VUgras{लोखंडी सामान} \textsuperscript{(67)} घडवतो, \VUgras{वीजकाम करणाऱ्यासाठी}
\textsuperscript{(124)} — जो \VUgras{गाडीखान्याच्या} \textsuperscript{(51)} छपरावर आहे.
त्याच्याजवळ \VUgras{भट्टी} \textsuperscript{(126)}, \VUgras{ऐरण} \textsuperscript{(127)} आणि
\VUgras{सांडशी} \textsuperscript{(128)} आहे.

%%K t05-01-42 p
वीजकाम करणारा दूरध्वनीची \VUgras{तांब्याची} \textsuperscript{(44)} \VUgras{तार} बसवतो आहे.

%%K t05-tit-7 sub
\VUpk{10.2pt}{40}{बागकाम.}
\VUsaut{3.0mm}

%%K t05-01-43 p
\textsc{काका}. --- \VUgras{अंगणाच्या} \textsuperscript{(45)} टोकाला, \VUgras{जाळीजवळ}
\textsuperscript{(47)} आणि \VUgras{फाटकाजवळ} \textsuperscript{(48)} तुला \VUgras{माळी}
\textsuperscript{(129)} दिसतो. तो \VUgras{लहान बाग} \textsuperscript{(49)} तयार करतो आहे.
आपल्या \VUgras{कुदळीने} \textsuperscript{(131)} त्याने जमीन खणली, बिया पेरल्या आणि
\VUgras{दंताळ्याने} \textsuperscript{(130)} सारखी केली. नंतर आपल्या \VUgras{झारीने}
\textsuperscript{(132)} तो \VUgras{वाफ्यांना} \textsuperscript{(150)} पाणी देईल, आणि रोपे वाढत
जातील.

%%K t05-01-44 p
\VUgras{मोतद्दाराने} \textsuperscript{(133)} नुकतीच घोड्याची निगा राखली आणि त्याला खाऊ घातले; आता तो \VUgras{गाडी} \textsuperscript{(52)} साफ करतो आहे — स्पंज, \VUgras{हातपंप} \textsuperscript{(134)} आणि बादलीभर पाणी घेऊन. नंतर तो ती पुन्हा गाडीखान्यात नेईल आणि जाड \VUgras{अडसराने} \textsuperscript{(53)} दार बंद करील.

%%K t05-01-45 p
आता तू खूश आहेस असे मला वाटते; तुला पुरेशी माहिती मिळाली.

%%K t05-01-46 p
\textsc{जॉन}. --- होय काका, पण मला घराचा आतला भागही आवडीने बघायचा आहे.

%%K t05-01-47 p
\textsc{काका}. --- ते उद्या होईल. श्री. रुफो तुला नकाशा समजावून सांगतील आणि प्रत्येक खोलीचे
नाव सांगतील.

%%K t05-tit-8 sub
\VUsaut{3.0mm}
\VUpk{10.2pt}{40}{घराची अंतर्गत रचना.}
\VUsaut{2.4mm}

%%K t05-apar-2 apar
{\centering\textit{(नकाशा पाहा.)}\par}
\VUsaut{2.4mm}

%%K t05-01-48 p
\textsc{वास्तुविशारद}. --- चल जॉन, मी तुला घराचे वेगवेगळे भाग लगेच दाखवतो.

%%K t05-01-49 p
आपण \VUgras{तळमजल्यात} \textsuperscript{(7)} शिरू. हे बघ \VUgras{घंटीचे}
\textsuperscript{(11)} बटण. आपण तीन \VUgras{पायऱ्या} \textsuperscript{(14)} चढतो —
\VUgras{चौथऱ्याच्या} \textsuperscript{(9)} —, \VUgras{उंबरठा} \textsuperscript{(10)} ओलांडतो —
\VUgras{प्रवेशद्वाराचा} \textsuperscript{(8)} — आणि आता आपण \VUgras{जिन्याच्या}
\textsuperscript{(13)} पायथ्याशी आहोत.

%%K t05-01-50 p
\textsc{जॉन}. --- ही मार्गिका आहे.

%%K t05-01-51 p
\textsc{वास्तुविशारद}. --- नाही, हा \VUgras{प्रवेशकक्ष} \textsuperscript{(12)} आहे.
\VUgras{मार्गिका} \textsuperscript{(a)} टोकाला आहे. उजवीकडे हे बघ \VUgras{दिवाणखाना}
\textsuperscript{(b)} आणि \VUgras{जेवणघर} \textsuperscript{(c)}; जेवणघरासमोर
\VUgras{स्वयंपाकघर} \textsuperscript{(d)} आणि \VUgras{नोकरांची खोली} \textsuperscript{(e)},
आणि त्यापुढे \VUgras{अभ्यासिका} \textsuperscript{(f)}. \VUgras{व्हरांडा}
\textsuperscript{(23)} आपल्या \VUgras{काचेच्या दारासह} \textsuperscript{(25)} दिवाणखान्याजवळ
आहे.

%%K t05-01-52 p
आता आपण जिना चढू. \VUgras{कठड्याचा} \textsuperscript{(15)} आधार घे आणि पायऱ्या मोज. त्या चौदा
आहेत. हा \VUgras{जिन्यावरचा चौक} \textsuperscript{(m)}; आपण \VUgras{पहिल्या मजल्यावर}
\textsuperscript{(27)} आहोत.

%%K t05-tit-9 sub
\VUsaut{3.0mm}
\VUpk{10.2pt}{40}{पहिला मजला.}
\VUsaut{3.0mm}

%%K t05-01-53 p
जिन्यासमोर \VUgras{शौचालय} \textsuperscript{(20)} आणि \VUgras{प्रसाधनगृह}
\textsuperscript{(j)} आहे. उजवीकडे सुंदर \VUgras{झोपायची खोली} \textsuperscript{(g)}, तिच्या
दोन खिडक्या घराच्या \VUgras{दर्शनी भागात} \textsuperscript{(1)}. डावीकडे \VUgras{न्हाणीघर}
\textsuperscript{(1)} आणि \VUgras{पाहुण्यांची खोली} \textsuperscript{(i)}, तिच्यात मोठा
\VUgras{कोनाडा} \textsuperscript{(h)}: झोपायच्या खोलीजवळ \VUgras{मुलांची खोली}
\textsuperscript{(k)} आहे. ती विस्तीर्ण \VUgras{गच्चीला} \textsuperscript{(21)} लागून आहे. या
गच्चीखाली दुसरे शौचालय (20) आहे, आणि भिंतीवर हे बघ जेवणासाठी वाजणारी \VUgras{घंटा}
\textsuperscript{(22)}.

%%K t05-01-54 p
\textsc{जॉन}. --- आपण \VUgras{दुसऱ्या मजल्यावर} जाऊ.

%%K t05-01-55 p
\textsc{वास्तुविशारद}. --- दुसरा मजला नाही. छपराखाली फक्त \VUgras{माड्या}
\textsuperscript{(33)} आणि कोठार (34) आहे. आपली फेरी संपली, आपण खाली उतरू शकतो.

%%K t05-01-56 p
\textsc{जॉन}. --- धन्यवाद महाशय, हे फार रोचक होते. मी जे पाहिले ते सगळे लगेच माझ्या वडिलांना
सांगणार आहे.

\VUsaut{4.0mm}
\VUfilet{22mm}
